%\begin{wrapfigure}{r}{0.6\textwidth}
% \begin{figure}[t]
%     \centering
%     \includegraphics[width=\linewidth]{figures/human_attention_inspired.pdf}
%     \caption{%Illustrations of I-JEPA and PEND.
%     \textbf{(A) Humans} perceive visual scenes selectively and sequentially, focusing on discriminative regions such as the red head, and red chest of a Northern Cardinal. \textbf{(B) \ours} emulates this process by ranking regions by attention-derived importance and predicting each next discriminative region's embedding in a sequential, GPT-style manner.}
%     \label{fig:human}
%     % \vspace{-0.5cm}
% \end{figure}
%\end{wrapfigure}

% \begin{figure}[h]
% {
%     \centering
%     \includegraphics[width=\textwidth]{figures/human_attention_inspired.pdf}
%     \setlength{\abovecaptionskip}{0.1cm}
%     \caption{
%     \textbf{(A) Humans} perceive visual scenes selectively and sequentially, focusing on discriminative regions such as the red head, and red chest. \textbf{(B) \ours} emulates this process by ranking regions by attention-derived importance and predicting each next discriminative region's embedding in a sequential manner.}
%     \label{fig:human}
% }
% \end{figure}
% In preamble


% In text
\begin{wrapfigure}{r}{0.57\textwidth}
    \centering
    % \vspace{-1.7\baselineskip} % optional: tighten top spacing
    \includegraphics[width=0.57\textwidth]{figures/human_attention_inspired.pdf}
    \setlength{\abovecaptionskip}{0.1cm}
    \caption{\textbf{(A) Humans} perceive visual scenes selectively and sequentially, focusing on discriminative regions.
    % such as the red head and red chest. 
    \textbf{(B) \ours} emulates this process by ranking regions by attention-derived importance and successively predicting next discriminative region embedding.}
    \label{fig:human}
    \vspace{-1.8\baselineskip} % optional: tighten bottom spacing
\end{wrapfigure}